\section{Finite support and the weights in the lifting
map}\label{finite-support-and-the-weights-in-the-lifting-map}

\subsection{FM0. Source, prerequisites and the question
calculated}\label{fm0.-source-prerequisites-and-the-question-calculated}

This calculation retains the finite adele coordinates in Alain Connes
and Caterina Consani's \href{https://arxiv.org/abs/1805.10501v1}{adelic
complex lift}, §5.2, equations \texttt{actionpq} and
\texttt{adelicomp3}. Their left action is \[
(X,Y)\longmapsto(aX+b,aY),\qquad a\in\mathbb Q^\times,
\ b\in\mathbb Q.
\tag{FM0.1}
\] The finite coordinates are part of that formula. They are not
numerical coordinates on the user's primitive \(Z_1/\tau\).

Operation prerequisites are the rational arithmetic already
reconstructed in the programme, finite-adele Schwartz--Bruhat tests,
their continuous linear dual, and the finite-dimensional primary blocks
TL2. These receiving vector spaces have addition. No sum involving
source \(\tau\) occurs. The current source definitions U01--U18 and the
complete global arguments USR-9ad1c0a2d09dba92, USR-f55d16feb948d8c2 and
USR-6152e3bc6302258c were consulted before this calculation. Their
locators and verbatim texts are retained in
\nolinkurl{CORPUS_AND_OPERATION_RULES.md} and
user\_corpus\_private/USER\_MESSAGES.jsonl. Complete arithmetic is used
as already reconstructed; no incomplete prime list substitutes for it.

The earlier archimedean uniqueness theorem DR6 is valid in its stated
receiving space. Here we compute precisely what changes after adjoining
particular finite-adele distributions. This is a comparison of proved
receiving constructions, not an assertion that the full tau
specialization has already been identified. The separate finite-boundary
sheaf calculation determines how these distributions enter the original
coefficient sheaves.

\subsection{FM1. All finite support profiles as actual
distributions}\label{fm1.-all-finite-support-profiles-as-actual-distributions}

For each rational prime \(p\), specify two numbers
\(\epsilon_{p,X},\epsilon_{p,Y}\in\{0,1\}\). They record which finite
coordinate is supported at zero. They are labels of the receiving
distribution, not the user's \(Z_2\) parity. Put \[
d_\epsilon(p)=\epsilon_{p,X}+\epsilon_{p,Y}\in\{0,1,2\}.
\tag{FM1.1}
\] Use the same additive Haar measure with
\(\operatorname{vol}(\mathbb Z_p)=1\) at every finite prime. Its role
and factors will remain explicit. In one coordinate write \(h_{p,0}\)
for the distribution given by Haar measure on \(\mathbb Q_p\), and
\(h_{p,1}=\delta_0\) for evaluation at zero. Both take value one on
\(\mathbf1_{\mathbb Z_p}\).

For a rational branch \(c\), translate the X factor by \(c\) when it is
a delta distribution; its Haar factor stays Haar. Define \[
h_{\epsilon,c}=\bigotimes_p
  \bigl(h_{p,\epsilon_{p,X},c}\otimes h_{p,\epsilon_{p,Y}}\bigr).
\tag{FM1.2}
\] This is a well-defined distribution on \(\mathbb A_f^2\), even when
the profile involves infinitely many primes. Here is a direct
construction. Every Schwartz--Bruhat test is a finite linear combination
of elementary tests whose p-components equal
\(\mathbf1_{\mathbb Z_p^2}\) outside a finite set. Enlarge that set to
include the denominator primes of \(c\). Define the value on an
elementary test as the finite product of its remaining local values.
Enlarging the set again does not change the product, because all extra
factors are one. This defines compatible linear functionals on the
finite tensor stages, hence a linear functional on their union. On each
compact-support, fixed-local-constancy stage the test space is finite
dimensional, so the functional is continuous there. The defining
inductive-limit topology of Schwartz--Bruhat tests therefore makes it a
distribution.

Equivalently it is Haar measure on the closed affine subgroup given by
\[
X_p=c\quad(\epsilon_{p,X}=1),\qquad
Y_p=0\quad(\epsilon_{p,Y}=1),
\tag{FM1.3}
\] embedded as a distribution in the ambient two-coordinate space. At
every unrestricted coordinate the Haar factor is retained. The subgroup
description follows by checking elementary tests against the product
definition. Its support is exactly (FM1.3): a test supported away from
that set is killed after a finite cylinder refinement; a neighbourhood
meeting that set contains a compact cylinder with positive product
measure on the unrestricted coordinates.

Distinct profiles at a fixed branch define linearly independent
distributions in every finite subfamily. To prove this, choose finitely
many primes which distinguish that subfamily. At each selected prime,
Haar measure and delta are independent: evaluation on
\(\mathbf1_{\mathbb Z_p^\times}\) is respectively \(1-p^{-1}\) and zero,
while evaluation on \(\mathbf1_{\mathbb Z_p}\) is one for both.
Translated versions give the same two-dimensional independence. Their
tensor products have independent coordinate functionals, by applying the
dual bases in each factor. Tests at the remaining primes take value one.
This separates every coefficient in the finite subfamily. Only algebraic
sums of profiles are used below.

\subsection{FM2. Every finite Jacobian and every prime
action}\label{fm2.-every-finite-jacobian-and-every-prime-action}

Let \(P_f(a,b)\) denote distribution pullback by (FM0.1) on
\(\mathbb A_f^2\). Our convention, for a finite-adele distribution
\(v\), is \[
\langle P_f(a,b)v,\varphi\rangle
=|a|_f^{-2}\langle v,\varphi((X-b)/a,Y/a)\rangle,
\quad |a|_f=\prod_p|a|_p.
\tag{FM2.1}
\] This is the pullback convention that gives ordinary composition for a
function regarded as a density distribution. Every product over absolute
values is finite because \(a\) is rational.

The exact transformation of (FM1.2) is \[
P_f(a,b)h_{\epsilon,c}
=\kappa_\epsilon(a)h_{\epsilon,(c-b)/a},
\quad
\kappa_\epsilon(a)=\prod_p|a|_p^{-d_\epsilon(p)}.
\tag{FM2.2}
\] Indeed a Haar coordinate supplies \(|a|_p\) on change of variable,
cancelling its ambient factor \(|a|_p^{-1}\). A delta coordinate
supplies evaluation without that change-of-variable factor, leaving
\(|a|_p^{-1}\). The X evaluation moves to \((c-b)/a\); the Y evaluation
stays at zero. Multiplication of these factors proves (FM2.2). In
particular \[
\kappa_\epsilon(p)=p^{d_\epsilon(p)}.
\tag{FM2.3}
\] For the profile with Haar in X and delta in Y at every finite prime,
the product formula gives \[
\kappa_\epsilon(a)=|a|_f^{-1}=|a|_\infty.
\tag{FM2.4}
\] This factor is lost if the finite coordinates are omitted. It is a
consequence of their actual distributional pullback, not a freely
selected twist.

\subsection{FM3. Combining finite support with the original primary
blocks}\label{fm3.-combining-finite-support-with-the-original-primary-blocks}

Fix an actual original-zeta zero \(\rho\), with multiplicity \(m\), and
retain \[
B_\rho=\mathbb C[x,y]/(x,y)^m,\qquad
J=M_x,\qquad q_\rho(x)=t,\quad q_\rho(y)=0.
\tag{FM3.1}
\] The j-th summand is \(y^j\mathbb C[x]/(x^{m-j})\), for \(0\le j<m\).
DR1 realizes its basis as \[
\Phi_c(x^uy^j)=\delta^{(j)}(X_\infty-c)\otimes
Y_\infty^{\rho+1}
\frac{(\log Y_\infty)^{m-j-1-u}}{(m-j-1-u)!},
\quad Y_\infty>0.
\tag{FM3.2}
\] Both archimedean coordinates are retained. The proof and original
differential are in DR1--DR5. Tensor (FM3.2) with (FM1.2). Tests and
distribution tensors exist as algebraic tensors of the already specified
continuous functionals; no product of arbitrary distributions is taken.

At \(c=0,b=0,a=p>0\), the exact eigenvalue and nilpotent action on this
block are \[
F_{p,\epsilon,j}
=p^{\rho-j+d_\epsilon(p)}
\exp((\log p)J_{m-j}).
\tag{FM3.3}
\] Here the archimedean factor is \(p^{\rho-j}\), including its
distribution Jacobian, and the finite factor is (FM2.3). Thus the
primewise absolute-value exponent is \[
2\Re\rho-2j+2d_\epsilon(p).
\tag{FM3.4}
\] We call this a primewise exponent; its identification with a Deligne
geometric weight requires the actual Frobenius and cohomology
comparison. In particular this formula does not assign a weight to
primitive \(\tau\).

\subsection{FM4. Exact classification of simultaneous prime
resonances}\label{fm4.-exact-classification-of-simultaneous-prime-resonances}

Take source zero \(\sigma\), source profile \(\epsilon^0\), and its full
primary block of multiplicity \(n=m_\sigma\). Take target zero \(\rho\),
target profile \(\epsilon^1\), and target normal order \(j\),
\(0\le j<m_\rho\). We calculate all linear maps between these
origin-branch blocks commuting with every rational prime action.

If the map is nonzero, the scalar eigenvalues of source and target must
agree at every prime. To justify this with the nilpotents retained, when
the two eigenvalues differ their annihilating powers \((X-\alpha)^n\)
and \((X-\beta)^{m_\rho-j}\) are coprime. Their polynomial Bezout
identity, applied across the intertwiner, forces the map to vanish.
Consequently a nonzero map requires \[
p^{\rho-\sigma-j+d_{\epsilon^1}(p)-d_{\epsilon^0}(p)}=1
\quad\text{for every prime }p.
\tag{FM4.1}
\] The original open critical strip gives \(-1<\Re\rho-\Re\sigma<1\).
Taking absolute values in (FM4.1) therefore forces \[
\Re\rho=\Re\sigma,\qquad
d_{\epsilon^1}(p)-d_{\epsilon^0}(p)=j
\quad\text{at every prime}.
\tag{FM4.2}
\] No finite prime sample replaces the second condition. Put
\(\rho-\sigma=i\theta\). At primes 2 and 3, (FM4.1) gives
\(\theta\log2=2\pi k_2\), \(\theta\log3=2\pi k_3\), with integers
\(k_2,k_3\). If \(\theta\ne0\), then \(2^{k_3}=3^{k_2}\) with nonzero
integers, contrary to unique factorization (clear negative powers by
multiplying). Hence \(\rho=\sigma\).

Conversely, when \[
\boxed{\rho=\sigma,\qquad
d_{\epsilon^1}(p)-d_{\epsilon^0}(p)=j\text{ for all }p,}
\tag{FM4.3}
\] the scalar factors coincide. Commuting with the remaining
exponentials is equivalent to \(J_{m-j}H=HT_m\). Indeed one may take the
finite polynomial logarithm of the unipotent matrices for the prime 2;
dividing by the nonzero real number \(\log2\) recovers their nilpotents.
This gives all maps explicitly: \[
H_v(f(t))=f(J_{m-j})v,
\quad v\in y^j\mathbb C[x]/(x^{m-j}).
\tag{FM4.4}
\] They are well-defined because \(J_{m-j}^m=0\), and every intertwiner
has this form by applying it to the cyclic generator. The dimension is
exactly \(m-j\). No spectral nilpotent has been dropped.

For a lift from a j=0 source into the kernel of \(q_\rho\), one has
\(j\ge1\). Since each finite profile has d between zero and two, (FM4.3)
restricts compensation to \(j=1\) or \(j=2\). At j=1 the difference must
be one at every prime; at j=2 it must be two at every prime. All higher
normal orders are excluded from this precise Haar/delta family by the
proved integer comparison. Other kinds of finite distributions are not
excluded by this calculation.

\subsection{FM5. A surviving lift variation, with the omitted finite
factor
restored}\label{fm5.-a-surviving-lift-variation-with-the-omitted-finite-factor-restored}

Let \(m\ge2\), let \(h_0\) be Haar in both finite coordinates, and let
\(h_1\) be Haar in finite X and delta in finite Y at every prime. Work
on the algebraic direct sum of their rational-branch blocks, retaining
both coefficients. Define a map into the kernel of q by \[
H(h_0\otimes f(t))=h_1\otimes y f(x),\qquad
H(h_1\otimes f(t))=0.
\tag{FM5.1}
\] The right side is taken modulo \((x,y)^m\), so the highest source
power is allowed to map to zero. It commutes with the full spectral
nilpotent \(J=M_x\) and every positive rational left dilation: the
finite factor \(a\) cancels exactly the normal-order factor \(a^{-1}\).
It commutes with rational translations because both distributions move
their rational branch by the same formula \((c-b)/a\). Thus, with the
earlier section i and quotient q on the two coefficient summands, \[
q(i+zH)=\mathrm{id}\quad(z\in\mathbb C).
\tag{FM5.2}
\] The sections are distinct since H sends the constant polynomial in
the h0 summand to the nonzero vector \(h_1\otimes y\). This proves
existence of a precisely computed family of lifts, not an obstruction to
existence.

No section here is asserted to intertwine the additional nilpotent
\(N=M_y\) with the zero operator on A. In fact \(Ni(1)=y\ne0\) for
m\textgreater1, and \(NH(1)=h_1\otimes y^2\ne0\) for m\textgreater2. The
stronger correct statement is about automorphisms of the full receiving
target: extend H to \(\widetilde H(h_0\otimes g)=h_1\otimes yg\),
\(\widetilde H(h_1\otimes g)=0\). Then \(\widetilde H^2=0\), and it
commutes with both J and N and every indicated positive affine action.
Thus \(U_z=I+z\widetilde H\) has inverse \(I-z\widetilde H\), satisfies
\(qU_z=q\), and sends i to the section \(i+zH\). Multiplication by y
proves commutation with both nilpotents directly; the finite and
archimedean factors in FM2--FM3 prove the affine identity.

The positive rational action is the one just proved. The full signed
action, arithmetic Frobenius, and all cochain factors are calculated
separately in \nolinkurl{FINITE_ADELE_FROBENIUS_DERIVATION.md}; no unsigned
formula is silently extended to negative a. The sheaf-local trace giving
the Haar-to-delta entry and the exact enlargement of coefficients are
calculated in \nolinkurl{FINITE_BOUNDARY_SHEAF_TRACE.md}. Those comparisons
determine the scope of (FM5.1) in the original geometry.

\subsection{FM6. Every external lattice label and the exact change in
the
argument}\label{fm6.-every-external-lattice-label-and-the-exact-change-in-the-argument}

For the earlier named receiving functor \(G_L\), lift every displayed
linear map f by \((v,\lambda)\mapsto(f(v),\lambda)\). This retains the
label at amplitude zero as well: it sends \((0,\lambda)\) to
\((0,\lambda)\). To verify addition and scalar action, insert the
component formulas in the already constructed receiving semimodules;
linearity gives the amplitude identity and both sides use the same join
or meet of labels. Composition follows by evaluation. These formulas are
not operations on primitive \(Z_1/\tau\).

The geometric support profile \(\epsilon\) is also retained. H changes
h0 to h1 and hence changes that geometric support. Preservation of an
external lattice label is not a claim that those geometric supports
coincide. Whether a required geometric filtration permits this map is
checked by the explicit sheaf trace, not inferred from the external
label formula.

Pierre Deligne's
\href{https://www.numdam.org/item/PMIHES_1980__52__137_0/}{Weil II},
§§3.6.1--3.6.3, proves vanishing after obtaining non-overlapping weight
ranges for the actual classes and actual support group. FM2--FM5 show
why the finite support factors must be calculated before importing that
step: in this specified part of the full adelic distribution space, they
restore exactly the eigenvalues that the archimedean-only computation
had separated. The original DR theorem remains true for its stated
archimedean receiver. Its uniqueness conclusion is not a theorem for
this larger coefficient space.

This gives a concrete remaining geometric contribution, with exact maps,
factors and nilpotents, to control in the requested tau lifting proof.
It neither supplies an RH counterexample nor establishes a contradiction
in the user's source. The original-zeta identification, source notation,
and current goal are unchanged.

\subsection{FM7. Quotienting the complete calculated ambiguity of
lifts}\label{fm7.-quotienting-the-complete-calculated-ambiguity-of-lifts}

We now retain every profile from FM1, every actual original zero and
every rational branch, taking algebraic direct sums so that each vector
has finitely many nonzero summands. Let \(\mathscr A_f\) and
\(\mathscr B_f\) be the sums of their A and B blocks, with q and i
defined blockwise. Each source profile is independently retained, even
when two profiles have the same numbers d(p). On \(\mathscr B_f\), J and
N denote the blockwise multiplications by x and y.

Consider all sections s of q that commute with J and with the complete
positive rational affine action just constructed. Their entire
difference from i is h=s−i, a map \(\mathscr A_f\to\ker q\). The
restrictions of h to an origin branch have finite-dimensional images,
hence finite support in profiles, zeros and branches. Equivariance under
dilation by 2 forces the finite set of output rational branches to be
stable under c↦c/2. This set can contain only zero: a nonzero rational
has an infinite orbit. Decompose that origin-branch image into its
independent profile and zero summands and its normal-order summands. FM4
proves that its only nonzero components are \[
h_{\epsilon^1,\epsilon^0,\rho}(f(t))
=y^j P_{\epsilon^1,\epsilon^0,\rho}(x) f(x),
\quad\deg P_{\epsilon^1,\epsilon^0,\rho}<m_\rho-j,
\tag{FM7.1}
\] where j is 1 or 2 and \(d_{\epsilon^1}(p)-d_{\epsilon^0}(p)=j\) for
every prime. All other components are zero. Translation equivariance
determines the same coefficients at every rational branch. Conversely
these formulas, with finitely many outputs per source summand, construct
every such h: FM4 gives all prime actions, the product formula extends
their scalar factors to every positive rational a, and the common branch
formula gives translations.

Extend each component of (FM7.1) to B by multiplication by the same
polynomial \(y^jP(x)\), and call the resulting block matrix
\(\widetilde h\). This is a well-defined operator on the algebraic
direct sum. It commutes with J, N and the full indicated affine action,
and \(q\widetilde h=0\). In a product of three such strictly raising
components, each finite prime's d-value would increase by at least
three. Since 0≤d(p)≤2, every such product is zero. Hence \[
\widetilde h^3=0,\qquad
U_h=I+\widetilde h,
\qquad U_h^{-1}=I-\widetilde h+\widetilde h^2.
\tag{FM7.2}
\] These are finite formulas even on the whole infinite profile sum.
Direct multiplication proves the inverse, while \(qU_h=q\) and
\(U_hi=i+h=s\).

Let \(\mathcal U\) be the group of all automorphisms \(I+T\) whose
entries have precisely this strictly raising form. Multiplication is \[
(I+T)(I+S)=I+(T+S+TS),
\tag{FM7.3}
\] and the profile increases and nilpotent polynomial products keep the
result in the same form. The inverse is FM7.2. This group acts
transitively on the calculated sections. It acts freely: an operator T
of this form is linear over the block algebra generated by x and y, and
the image of i contains the generator 1 of each source B block. If Ti=0,
every column is zero on that generator, so every polynomial entry and T
itself vanish. Therefore \[
\boxed{\text{All calculated prime-and-J-equivariant lifts form one free orbit of }\mathcal U.}
\tag{FM7.4}
\] Freeness at i implies freeness at every other section: transitivity
expresses that section as Ui, and conjugation by U identifies its
stabilizer with the stabilizer of i.

This is an exact quotient of the full ambiguity calculated here,
including both nilpotents in the automorphisms, every finite support
profile, all primes and all rational branches. It realizes a part of the
user's global-quotient proposal as a proved map and action, without
assuming that every conceivable full-geometric contribution belongs to
this Haar/delta family. The original zero character \(p^\rho\), its
multiplicity and its nilpotent are unchanged under this quotient. The
construction therefore removes this lifting ambiguity but does not force
\(\Re\rho=1/2\). Its extension to the full geometric specialization
remains the active goal.
